\scanpage{59}

\begin{center}- 154 -\end{center}
\begin{flushright}IX\end{flushright}

de Lie, Chap IV, Annexe) $\Gamma(C)$ dont les sommets sont les composantes irréductibles de $\widetilde C$, deux sommets $\widetilde C_i$ et $\widetilde C_j$ étant reliés par un ensemble d'arrêtes en correspondance bijective avec l'ensemble $\widetilde C_i\times_C\widetilde C_j$ des ensembles de deux branches \uline{distinctes} de $C$ contenues dans $\widetilde C_i$ resp. $\widetilde C_j$ passant par le même point de $C$, ensemble égal à $\widetilde C_i\times_C\widetilde C_j$ si $i\neq j$, à $\widetilde C_i\times_C\widetilde C_i-\Delta(\widetilde C_i)$ modulo symétrie $(\Delta:\widetilde C\longrightarrow\widetilde C\times_C\widetilde C$ étant la diagonale) si $i=j$ (donc vide si $C_i$ est géométriquement unibranche). Ceci posé, on a des isomorphismes canoniques
\[
(12.3.7)\qquad M=D(T)\xleftarrow{\sim}H_1(\Gamma(C),\mathbb Z),
\quad\text{donc }\check M=R\simeq H^1(\Gamma(C),\mathbb Z)\quad.
\]
Nous n'aurons pas besoin de cette interprétation combinatoire de $M$, et allons nous borner à définir l'homomorphisme canonique
\[
H_1(\Gamma(C),\mathbb Z)\longrightarrow M\quad,
\]
laissant au lecteur le soin de prouver que c'est un isomorphisme. Soit donc $I$ l'ensemble des points de $C$ appartenant a au moins deux branches de $C$, $J$ (comme ci-dessus) l'ensemble des composantes irréductibles de $\widetilde C$ (ou de $C$, cela revient au même), $K$ l'ensemble des ``branches'' de $C$ en les points de $I$, i.e. l'image inverse de $I$ dans $\widetilde C$. On a donc une application canonique
\[
(12.3.9)\qquad h:K\longrightarrow I\times J,\quad h=(f,g)\quad,
\]
et la description (12.3.4) de $M$ s'exprime en termes de $h$ comme
\[
(12.3.10)\qquad M\simeq\operatorname{Ker}(\mathbb Z^{(h)}:\mathbb Z^{(K)}
\longrightarrow\mathbb Z^{(I\times J)})\quad,
\]
car l'expression (12.3.3) de $R=\check M$ peut s'interpréter manifestement

\margin{472}
\margin{472}
\margin{58}
